\documentclass[aspectratio=169,11pt]{beamer}

% Самостоятельный исходник: pdfLaTeX, без внешних рисунков и файлов.
\usepackage{iftex}
\ifPDFTeX
  \usepackage[T2A]{fontenc}
  \usepackage[utf8]{inputenc}
\else
  \usepackage{fontspec}
  \setmainfont{cmunrm.otf}[BoldFont=cmunbx.otf,ItalicFont=cmunti.otf,BoldItalicFont=cmunbi.otf]
  \setsansfont{cmunss.otf}[BoldFont=cmunsx.otf,ItalicFont=cmunsi.otf]
  \setmonofont{cmuntt.otf}
\fi
\usepackage[russian]{babel}
\usepackage{amsmath,amssymb,amsfonts,array}
\usepackage{tikz}
\usetikzlibrary{arrows.meta}

\usetheme{Madrid}
\usefonttheme{serif}
\setbeamercolor{background canvas}{bg=white}
\setbeamercolor{structure}{fg=blue!70!black}
\setbeamertemplate{navigation symbols}{}
\setbeamertemplate{footline}{%
  \leavevmode\hbox{%
  \begin{beamercolorbox}[wd=.28\paperwidth,ht=2.6ex,dp=1.2ex,center]{author in head/foot}%
  \insertshortauthor{} (\insertshortinstitute)\end{beamercolorbox}%
  \begin{beamercolorbox}[wd=.50\paperwidth,ht=2.6ex,dp=1.2ex,center]{title in head/foot}%
  Обыкновенные дифференциальные уравнения\end{beamercolorbox}%
  \begin{beamercolorbox}[wd=.22\paperwidth,ht=2.6ex,dp=1.2ex,center]{date in head/foot}%
  Лекция 2 \quad \insertframenumber/\inserttotalframenumber\end{beamercolorbox}}}
\newcommand{\R}{\mathbb R}
\newcommand{\dd}{\,d}
\title[Простейшие методы]{Обыкновенные дифференциальные уравнения}
\subtitle{Лекция 2. Простейшие методы решения\\Автономные уравнения и метод Эйлера}
% Лектор и место проведения курса.
\author[Ховрычев А. А.]{Ховрычев Андрей Александрович}
\institute[УрФУ]{ФИИТ, УрФУ}
\date{}


\begin{document}
% 1
\begin{frame}
\titlepage
\centering Формула, фазовая прямая и численное приближение
\end{frame}

% 2
\begin{frame}{От проверки решения к методам его нахождения}
На прошлой лекции: решение задаётся функцией \emph{на интервале},
а задача Коши добавляет начальные данные.
\begin{block}{Сегодня}
\begin{enumerate}
\item Разделение переменных: получить явную или неявную формулу.
\item Фазовая прямая: понять движение без интегрирования.
\item Интегрирующий множитель: решить линейное уравнение.
\item Метод Эйлера: построить приближение и оценить его ограничения.
\end{enumerate}
\end{block}
В примерах проверяем не только равенство, но и область определения,
постоянные решения, начальное условие и поведение при больших временах.
\end{frame}

\section{Разделение переменных}
% 3
\begin{frame}{Самый простой случай: $x'=g(t)$}
Пусть $g$ непрерывна на интервале $J$ и $t_0\in J$.
Тогда задача
\[
x'=g(t),\qquad x(t_0)=x_0
\]
имеет решение
\[
\boxed{x(t)=x_0+\int_{t_0}^{t}g(s)\,ds},\qquad t\in J.
\]
По формуле Ньютона--Лейбница производная равна $g(t)$.
Если $x_1,x_2$ --- два решения, то $(x_1-x_2)'=0$;
при одинаковых начальных данных разность равна нулю.
\begin{exampleblock}{Пример}
$x'=2t$, $x(1)=0$ дают $x(t)=t^2-1$.
\end{exampleblock}
\end{frame}

% 4
\begin{frame}{Уравнение с разделяющимися переменными}
\begin{block}{Определение}
Уравнение допускает \alert{разделение переменных}, если имеет вид
\[
x'=g(t)q(x).
\]
Пусть $g$ непрерывна на $J$, а $q$ непрерывна и не равна нулю
на интервале $U$, содержащем $x_0$.
\end{block}
Для решения с $x(t_0)=x_0$, пока $x(t)\in U$, получаем
\[
\boxed{\int_{x_0}^{x(t)}\frac{du}{q(u)}
=\int_{t_0}^{t}g(s)\,ds}.
\]
Нули $q$ требуют отдельного рассмотрения: делить на $q(x)$ там нельзя.
\end{frame}

% 5
\begin{frame}{Почему разделение переменных работает?}
Положим
\[
H(x)=\int_{x_0}^x\frac{du}{q(u)},\qquad
G(t)=\int_{t_0}^t g(s)\,ds.
\]
Вдоль решения по цепному правилу
\[
\frac{d}{dt}H(x(t))=\frac{x'(t)}{q(x(t))}=g(t)=G'(t).
\]
Значит, $H(x(t))=G(t)$: константа равна нулю при $t=t_0$.
\begin{block}{Обратный переход}
$H'(x)=1/q(x)\ne0$, поэтому $H$ строго монотонна на $U$
и имеет дифференцируемую обратную.
Функция $x(t)=H^{-1}(G(t))$ решает задачу там, где $G(t)\in H(U)$.
\end{block}
\end{frame}

% 6
\begin{frame}{Разделение переменных: порядок действий}
\begin{enumerate}
\item Найти нули $q(a)=0$ и записать постоянные решения $x(t)\equiv a$.
\item На каждом интервале $q(x)\ne0$ интегрировать:
\[
\int\frac{dx}{q(x)}=\int g(t)\,dt+C.
\]
\item Использовать начальные данные, выбрать нужную ветвь.
\item Указать интервал решения и проверить результат подстановкой.
\end{enumerate}
\begin{alertblock}{Логическая оговорка}
Без единственности возможны склейки постоянных и непостоянных частей.
Список «постоянные решения плюс формулы на $q\ne0$»
не всегда исчерпывает все решения.
\end{alertblock}
\end{frame}

% 7
\begin{frame}{Пример: экспонента с переменным темпом}
\[
x'=2tx,\qquad x(0)=x_0.
\]
На участке $x\ne0$:
\[
\frac{dx}{x}=2t\,dt,\qquad \ln|x|=t^2+C,
\qquad x=Ae^{t^2},\quad A\ne0.
\]
При $x_0=0$ отдельно получаем $x\equiv0$.
В итоге
\[
\boxed{x(t)=x_0e^{t^2}},\qquad t\in\R.
\]
\begin{block}{Проверка полноты}
$\bigl(e^{-t^2}x(t)\bigr)'=0$ для любого решения.
Поэтому других решений нет, включая случай $x_0=0$.
\end{block}
\end{frame}

% 8
\begin{frame}{Пример: формула и её максимальный интервал}
\[
x'=x^2,\qquad x(0)=x_0.
\]
При $x_0\ne0$ разделение даёт
\[
-\frac1x=t-\frac1{x_0},\qquad
\boxed{x(t)=\frac{x_0}{1-x_0t}}.
\]
\[
I_{\max}=
\begin{cases}
(-\infty,1/x_0),&x_0>0,\\
(1/x_0,+\infty),&x_0<0.
\end{cases}
\]
При $x_0=0$ решение $x\equiv0$ определено на всей $\R$.
\begin{alertblock}{Выбор интервала}
Берём компоненту области формулы, содержащую начальный момент $0$.
В полюсе конечное непрерывное продолжение невозможно.
\end{alertblock}
\end{frame}

% 9
\begin{frame}{Неявная формула тоже описывает решение}
\[
x'=\frac{t}{1+x^2},\qquad x(0)=0.
\]
Умножая на $1+x^2$ и интегрируя, получаем
\[
\boxed{x+\frac{x^3}{3}=\frac{t^2}{2}}.
\]
Функция $H(x)=x+x^3/3$ строго возрастает, поскольку $H'=1+x^2>0$,
и отображает $\R$ на $\R$.
\begin{block}{Что можно заключить без формулы для корня кубического уравнения?}
Для каждого $t\in\R$ значение $x(t)$ определяется однозначно.
Решение существует на всей $\R$, чётно, неотрицательно
и возрастает при $t>0$.
\end{block}
\end{frame}

% 10
\begin{frame}{Деление может скрыть часть движения}
Для $x'=2\sqrt{|x|}$, $x(0)=0$ из первой лекции известны решения
\[
x_a(t)=
\begin{cases}
0,&t\le a,\\
(t-a)^2,&t>a,
\end{cases}
\qquad a\ge0.
\]
Разделение на участке $x>0$ даёт $\sqrt{x}=t-a$,
то есть только часть при $t>a$.
\begin{alertblock}{Условие важно}
Если $q$ лишь непрерывна, решение может прийти в нуль $q$
или покинуть его. При $q\in C^1$ локальная единственность
исключает такое поведение; доказательство --- в лекции 3.
\end{alertblock}
\end{frame}

\section{Автономные уравнения на прямой}
% 11
\begin{frame}{Автономность: закон не зависит от часов}
Рассмотрим $x'=q(x)$; в этом разделе предполагаем $q\in C^1$.
\begin{itemize}
\item Нули $q$ задают равновесия.
\item На интервалах $q>0$ решения возрастают, на интервалах $q<0$ убывают.
\item Непостоянное решение не проходит через равновесие
при выполнении локальной единственности.
\end{itemize}
\begin{block}{Сдвиг времени}
Если $\varphi(t)$ --- решение на $I$, то $\psi(t)=\varphi(t+c)$
--- решение на $I-c$:
\[
\psi'(t)=q(\varphi(t+c))=q(\psi(t)).
\]
\end{block}
Здесь используем локальную единственность для $C^1$-поля;
её обоснование будет дано в следующей лекции.
\end{frame}

% 12
\begin{frame}{Как читать фазовую прямую}
\begin{enumerate}
\item Отметить нули $q$.
\item Определить знак $q$ между соседними нулями.
\item Поставить стрелки: вправо при $q>0$, влево при $q<0$.
\end{enumerate}
\begin{block}{О пределе решения}
Если решение существует при всех $t\ge0$ и $x(t)\to\ell\in\R$,
где $\ell$ лежит внутри области $q$, то $q(\ell)=0$.
\end{block}
Действительно, при $q(\ell)>0$ вблизи $\ell$ было бы $x'\ge c>0$,
что несовместимо с конечным пределом. Случай $q(\ell)<0$ аналогичен.
\medskip

Для решения, запертого в компактном отрезке области $C^1$-поля,
продолжимость вперёд будет обоснована в лекции 5.
\end{frame}

% 13
\begin{frame}{Поведение вблизи изолированного равновесия}
\begin{center}
\begin{tikzpicture}[>=Stealth,x=1cm,y=1cm]
\foreach \y in {0,1.1,2.2}{
\draw[gray] (-2,\y)--(2,\y);
\fill (0,\y) circle (2pt);
}
\draw[blue!75!black,thick,->] (-1.8,2.2)--(-.4,2.2);
\draw[blue!75!black,thick,->] (1.8,2.2)--(.4,2.2);
\node[right] at (2.3,2.2) {Притягивает с обеих сторон};
\draw[red!75!black,thick,->] (-.4,1.1)--(-1.8,1.1);
\draw[red!75!black,thick,->] (.4,1.1)--(1.8,1.1);
\node[right] at (2.3,1.1) {Отталкивает с обеих сторон};
\draw[teal!75!black,thick,->] (-1.8,0)--(-.4,0);
\draw[teal!75!black,thick,->] (.4,0)--(1.8,0);
\node[right] at (2.3,0) {Притягивает только слева};
\end{tikzpicture}
\end{center}
Если $q(a)=0$, то $q'(a)<0$ даёт первый случай,
$q'(a)>0$ --- второй. При $q'(a)=0$ нужно исследовать знаки.
\begin{block}{Терминология}
Первый случай --- асимптотически устойчивое равновесие;
второй --- неустойчивое. Третий называют полуустойчивым,
но на всей прямой оно неустойчиво по Ляпунову.
\end{block}
\end{frame}

% 14
\begin{frame}{Пример: три равновесия}
\[
x'=x(1-x^2),\qquad q'(x)=1-3x^2.
\]
\begin{center}
\begin{tikzpicture}[>=Stealth,x=2cm]
\draw[->] (-2,0)--(2.1,0) node[right] {$x$};
\foreach \a in {-1,0,1} \fill (\a,0) circle (2pt) node[below=4pt] {$\a$};
\draw[blue!75!black,very thick,->] (-1.85,0)--(-1.2,0);
\draw[blue!75!black,very thick,->] (-.15,0)--(-.8,0);
\draw[blue!75!black,very thick,->] (.15,0)--(.8,0);
\draw[blue!75!black,very thick,->] (1.85,0)--(1.2,0);
\foreach \x/\s in {-1.5/+, -.5/-, .5/+, 1.5/-} \node at (\x,.4) {$\s$};
\end{tikzpicture}
\end{center}
\begin{itemize}
\item $-1$ и $1$ притягивают; $0$ отталкивает.
\item При $x_0>0$ решение стремится к $1$, при $x_0<0$ --- к $-1$.
\item При $x_0=0$ решение остаётся в нуле.
\end{itemize}
Непостоянные решения здесь монотонны: периодических колебаний нет.
\end{frame}

% 15
\begin{frame}{Логистическая модель: ограниченный рост}
Экспоненциальная модель $x'=rx$ предполагает постоянный
темп роста на одну особь. Учтём ограниченность ресурсов:
\[
\frac{x'}x=r\left(1-\frac{x}{K}\right),\qquad x>0.
\]
Получаем
\[
\boxed{x'=rx\left(1-\frac{x}{K}\right)},\qquad r>0,\quad K>0.
\]
\begin{itemize}
\item $r$ --- темп роста при малой численности, размерность $1/\text{время}$.
\item $K$ --- ёмкость среды, размерность численности.
\item Биологическая область: $x\ge0$; $x=0$ проверяем отдельно.
\end{itemize}
Это непрерывная модель численности с неизменными параметрами среды.
\end{frame}

% 16
\begin{frame}{Логистическая модель: сначала геометрия}
\begin{center}
\begin{tikzpicture}[>=Stealth,x=1.5cm]
\draw[->] (-1.6,0)--(4.8,0) node[right] {$x$};
\fill (0,0) circle (2pt) node[below] {$0$};
\fill (3,0) circle (2pt) node[below] {$K$};
\draw[gray,thick,->] (-.2,0)--(-1.3,0);
\draw[blue!75!black,very thick,->] (.4,0)--(2.6,0);
\draw[blue!75!black,very thick,->] (4.4,0)--(3.4,0);
\node at (-.8,.4) {$-$}; \node at (1.5,.4) {$+$}; \node at (3.9,.4) {$-$};
\end{tikzpicture}
\end{center}
\begin{itemize}
\item $0<x_0<K$: численность растёт, оставаясь ниже $K$.
\item $x_0>K$: численность убывает, оставаясь выше $K$.
\item $x_0=0$ и $x_0=K$: постоянные решения.
\end{itemize}
\begin{block}{Прогноз}
При любом $x_0>0$ ожидаем $x(t)\to K$ при $t\to+\infty$.
Теперь подтвердим и глобальность вперёд, и предел явной формулой.
\end{block}
\end{frame}

% 17
\begin{frame}{Логистическое уравнение: разделяем переменные}
На участках $x\ne0,K$:
\[
\frac{dx}{x(1-x/K)}=r\,dt,\qquad
\frac{1}{x(1-x/K)}=\frac1x+\frac1{K-x}.
\]
Поэтому
\[
\ln|x|-\ln|K-x|=rt+C,\qquad
\frac{x}{K-x}=Be^{rt}.
\]
После подстановки $x(0)=x_0>0$ и разрешения относительно $x$:
\[
\boxed{x(t)=\frac{K}{1+\left(K/x_0-1\right)e^{-rt}}}.
\]
Полученная формула также включает $x_0=K$ после отдельной проверки.
Решение $x\equiv0$ записываем отдельно.
\end{frame}

% 18
\begin{frame}{Логистическая формула: область и предел}
Обозначим $A=K/x_0-1$, где $x_0>0$.
\[
x(t)=\frac{K}{1+Ae^{-rt}}.
\]
\begin{itemize}
\item Если $0<x_0<K$, то $A>0$: знаменатель положителен на всей $\R$.
\item Если $x_0=K$, то $A=0$: решение постоянно.
\item Если $x_0>K$, то $-1<A<0$: для $t\ge0$ знаменатель
не меньше $1+A=K/x_0>0$.
\end{itemize}
Во всех трёх случаях решение определено при всех $t\ge0$ и
\[
\lim_{t\to+\infty}x(t)=K.
\]
\alert{В прошлое} при $x_0>K$ есть полюс:
$t_*=\frac1r\ln(1-K/x_0)<0$; максимальный интервал --- $(t_*,+\infty)$.
\end{frame}

% 19
\begin{frame}{Семейство логистических кривых}
\begin{columns}[c]
\begin{column}{.59\textwidth}
\centering
\begin{tikzpicture}[x=1.05cm,y=2.05cm,>=Stealth]
\draw[->] (-.2,0)--(5.2,0) node[right] {$rt$};
\draw[->] (0,-.1)--(0,1.95) node[above] {$x/K$};
\draw[dashed,gray] (0,1)--(4.8,1) node[right] {$1$};
\draw[dotted,gray] (0,.5)--(4.8,.5);
\node[left] at (0,.5) {$1/2$};
\foreach \a in {.1,.3,.7,1.7}{
\draw[blue!70!black,thick,domain=0:4.8,samples=90]
plot (\x,{1/(1+(1/\a-1)*exp(-\x))});
}
\draw[teal!70!black,thick] (0,1)--(4.8,1);
\end{tikzpicture}
\end{column}
\begin{column}{.39\textwidth}
При $0<x<K$:
\[
x''=r(1-2x/K)x'.
\]
Скорость роста максимальна при $x=K/2$:
\[
x'_{\max}=\frac{rK}{4}.
\]
Если $x_0<K/2$, впереди есть точка перегиба.
Если $x_0>K$, решение убывает к $K$.
\end{column}
\end{columns}
\end{frame}

% 20
\begin{frame}{Когда популяция достигнет заданного уровня?}
Пусть $x_0=K/4$. Тогда
\[
x(t)=\frac{K}{1+3e^{-rt}}.
\]
Найдём время достижения $3K/4$:
\[
\frac34=\frac{1}{1+3e^{-rt}}
\quad\Longleftrightarrow\quad e^{-rt}=\frac19,
\qquad \boxed{t=\frac{\ln9}{r}}.
\]
\begin{block}{Интерпретация}
Время обратно пропорционально $r$.
Параметр $K$ сократился, поскольку заданы доли ёмкости среды.
\end{block}
Уровень $K$ достигается лишь в пределе, если $x_0\ne K$.
\end{frame}

\section{Линейные уравнения первого порядка}
% 21
\begin{frame}{Линейность --- другой путь к явному решению}
\begin{block}{Определение}
Линейное уравнение первого порядка:
\[
\boxed{x'+a(t)x=b(t)},
\]
где $a,b$ непрерывны на интервале $J$.
При $b\equiv0$ оно называется однородным.
\end{block}
Неизвестная и её производная входят линейно;
коэффициенты могут зависеть от времени.
\[
x'+tx=\sin t\quad\text{--- линейное},\qquad
x'=x^2\quad\text{--- нелинейное}.
\]
Не всякое линейное уравнение разделяется на множители $g(t)q(x)$.
\end{frame}

% 22
\begin{frame}{Как найти интегрирующий множитель?}
Хотим превратить левую часть в производную произведения:
\[
(\mu x)'=\mu x'+\mu'x=\mu(x'+ax).
\]
Для этого достаточно $\mu'=a\mu$. Возьмём
\[
\mu(t)=\exp\left(\int_{t_0}^t a(s)\,ds\right)>0,\qquad \mu(t_0)=1.
\]
Умножая уравнение на $\mu$, получаем
\[
\boxed{(\mu x)'=\mu b}.
\]
\begin{block}{Почему преобразование эквивалентно?}
$\mu$ нигде не обращается в нуль, поэтому можно и умножить на неё,
и разделить обратно, не теряя решений.
\end{block}
\end{frame}

% 23
\begin{frame}{Формула решения линейной задачи Коши}
Интегрируя $(\mu x)'=\mu b$ от $t_0$ до $t$, получаем
\[
\mu(t)x(t)-x_0=\int_{t_0}^t\mu(s)b(s)\,ds.
\]
Следовательно,
\[
\boxed{x(t)=\frac1{\mu(t)}
\left[x_0+\int_{t_0}^t\mu(s)b(s)\,ds\right]},\qquad t\in J.
\]
\begin{block}{Существование и единственность уже доказаны}
Правая часть определена и дифференцируема на всём $J$,
удовлетворяет уравнению и начальному условию.
Любое решение обязано иметь эту форму.
\end{block}
Интегралы не обязаны выражаться через элементарные функции.
\end{frame}

% 24
\begin{frame}{Возвращаемся к полю $x'=t-x$}
\[
x'+x=t,\qquad x(0)=x_0,\qquad \mu(t)=e^t.
\]
Тогда
\[
(e^tx)'=te^t,\qquad
\int_0^t se^s\,ds=(t-1)e^t+1.
\]
Отсюда
\[
\boxed{x(t)=t-1+(x_0+1)e^{-t}},\qquad t\in\R.
\]
\begin{itemize}
\item $t-1$ --- частное решение неоднородного уравнения.
\item $(x_0+1)e^{-t}$ --- затухающее решение однородного уравнения.
\item Разность любых двух решений имеет вид $Ce^{-t}$.
\end{itemize}
Это семейство мы видели на рисунке в первой лекции.
\end{frame}

% 25
\begin{frame}{Как выбрать метод?}
\renewcommand{\arraystretch}{1.5}
\begin{tabular}{>{\raggedright\arraybackslash}p{.28\textwidth}
>{\raggedright\arraybackslash}p{.29\textwidth}
>{\raggedright\arraybackslash}p{.33\textwidth}}
\textbf{Вид} & \textbf{Первый шаг} & \textbf{Что проверить} \\
$x'=g(t)$ & Интегрирование & Начальное условие \\
$x'=g(t)q(x)$ & Разделение & Нули $q$, ветви, интервал \\
$x'=q(x)$ & Фазовая прямая & Равновесия, знаки, продолжимость \\
$x'+a(t)x=b(t)$ & Множитель $\mu$ & Непрерывность коэффициентов \\
Общее $x'=f(t,x)$ & Численное приближение & Шаг и ошибка
\end{tabular}
\medskip

Одно уравнение может относиться к нескольким строкам.
Например, $x'=-kx$ и линейно, и автономно, и разделяется.
\end{frame}

\section{Метод Эйлера}
% 26
\begin{frame}{От касательной к пошаговому алгоритму}
Для задачи $x'=f(t,x)$, $x(t_0)=x_0$ выберем шаг $h>0$.
Пусть $t_n=t_0+nh$, а $X_n$ --- приближение к $x(t_n)$.
\[
\boxed{X_0=x_0,\qquad X_{n+1}=X_n+h f(t_n,X_n)}.
\]
\begin{enumerate}
\item Вычислить наклон в текущей приближённой точке $(t_n,X_n)$.
\item Пройти время $h$ вдоль соответствующей прямой.
\item Повторить с новым значением $X_{n+1}$.
\end{enumerate}
\begin{alertblock}{Различаем обозначения}
$x(t_n)$ --- точное значение; $X_n$ --- вычисленное.
Наклон на следующем шаге берётся в приближённой точке.
\end{alertblock}
\end{frame}

% 27
\begin{frame}{Ломаная Эйлера для $x'=-x$, $x(0)=1$}
\begin{columns}[c]
\begin{column}{.61\textwidth}
\centering
\begin{tikzpicture}[x=2cm,y=2.6cm,>=Stealth]
\draw[->] (-.12,0)--(2.5,0) node[right] {$t$};
\draw[->] (0,-.08)--(0,1.2) node[above] {$x$};
\draw[blue!75!black,very thick,domain=0:2.25,samples=70] plot (\x,{exp(-\x)});
\draw[red!75!black,thick] (0,1)--(.5,.5)--(1,.25)--(1.5,.125)--(2,.0625);
\foreach \t/\y in {0/1,.5/.5,1/.25,1.5/.125,2/.0625}{
\fill[red!75!black] (\t,\y) circle (1.5pt);
}
\foreach \t in {.5,1,1.5,2} \node[below] at (\t,0) {\scriptsize $\t$};
\node[left] at (0,1) {$1$};
\end{tikzpicture}
\end{column}
\begin{column}{.37\textwidth}
При $h=1/2$:
\[
X_{n+1}=\tfrac12X_n,\qquad X_n=2^{-n}.
\]
Синяя кривая: $e^{-t}$.

Красная ломаная: метод Эйлера.
\medskip

У ломаной обычно есть изломы; она не является классическим решением ОДУ.
\end{column}
\end{columns}
\end{frame}

% 28
\begin{frame}{Численный эксперимент: уменьшаем шаг}
Для $x'=-x$, $x(0)=1$, в точке $T=1$:
\[
X_N=(1-h)^N,\qquad Nh=1,\qquad x(1)=e^{-1}\approx0.367879.
\]
\begin{center}
\renewcommand{\arraystretch}{1.25}
\begin{tabular}{c r r r}
$h$ & $N$ & $X_N$ & $|X_N-x(1)|$\\ \hline
$0.2$ & 5 & $0.327680$ & $0.040199$\\
$0.1$ & 10 & $0.348678$ & $0.019201$\\
$0.05$ & 20 & $0.358486$ & $0.009394$\\
$0.025$ & 40 & $0.363232$ & $0.004647$
\end{tabular}
\end{center}
При делении шага пополам ошибка здесь уменьшается примерно вдвое.
\begin{block}{Наблюдение и теорема}
Таблица подсказывает порядок $O(h)$, но один эксперимент не доказывает его
для произвольного уравнения.
\end{block}
\end{frame}

% 29
\begin{frame}{Ошибка одного шага и накопленная ошибка}
Пусть точное решение дважды дифференцируемо и $|x''|\le M$
на фиксированном отрезке $[t_0,T]$.
Остаток шага, начатого с \emph{точного} значения:
\[
\delta_n=x(t_{n+1})-x(t_n)-h f(t_n,x(t_n)),\qquad
|\delta_n|\le\frac{M}{2}h^2.
\]
\begin{block}{Два разных порядка}
Ошибка одного шага $\delta_n$ имеет порядок $O(h^2)$.
Глобальная ошибка $e_n=x(t_n)-X_n$ при условии Липшица по $x$
на общей области имеет порядок $O(h)$ на фиксированном отрезке.
\end{block}
Шагов порядка $1/h$; ошибки переносятся и накапливаются.
Оценка относится к точной арифметике и требует оставаться в области,
где действуют гипотезы.
\end{frame}

% 30
\begin{frame}{Дополнение: от локальной ошибки к глобальной}
\small
Пусть на $[t_0,T]\times\R$ выполнено
$|f(t,u)-f(t,v)|\le L|u-v|$, а $|x''|\le M$ вдоль точного решения.
Тогда
\[
e_{n+1}=e_n+h\bigl(f(t_n,x(t_n))-f(t_n,X_n)\bigr)+\delta_n,
\]
\[
|e_{n+1}|\le(1+hL)|e_n|+\frac M2h^2,\qquad e_0=0.
\]
При $L>0$ сумма геометрической прогрессии даёт
\[
|e_n|\le\frac{Mh}{2L}\bigl((1+hL)^n-1\bigr)
\le\boxed{\frac{Mh}{2L}\bigl(e^{L(T-t_0)}-1\bigr)}.
\]
При $L=0$: $|e_n|\le Mh(T-t_0)/2$.
\begin{alertblock}{Граница вывода}
$T$ фиксировано. Эта оценка не обещает точности при $t\to+\infty$
для фиксированного шага. Доказательство можно разобрать на практике.
\end{alertblock}
\end{frame}

% 31
\begin{frame}{Слишком большой шаг меняет качественную картину}
Для $x'=-kx$, $k>0$, точное решение $x_0e^{-kt}$ затухает.
Метод Эйлера даёт
\[
X_n=(1-hk)^n x_0.
\]
При $x_0\ne0$:
\[
\boxed{X_n\to0\quad\Longleftrightarrow\quad |1-hk|<1
\quad\Longleftrightarrow\quad 0<hk<2}.
\]
\begin{itemize}
\item $0<hk<1$: знак сохраняется, модуль убывает.
\item $hk=1$: после одного шага все значения равны нулю.
\item $1<hk<2$: затухающие колебания знака.
\item $hk=2$: незатухающие колебания; $hk>2$: рост модуля.
\end{itemize}
\alert{Затухание сеточного решения ещё не гарантирует хорошую точность.}
\end{frame}

% 32
\begin{frame}{Одно уравнение --- разное поведение сетки}
\begin{columns}[c]
\begin{column}{.60\textwidth}
\centering
\begin{tikzpicture}[x=.6cm,y=1.2cm,>=Stealth]
\draw[->] (-.2,0)--(8.6,0) node[right] {$t$};
\draw[->] (0,-1.7)--(0,1.85) node[above] {$x$};
\draw[blue!75!black,very thick,domain=0:8,samples=70] plot (\x,{exp(-\x)});
\draw[teal!75!black,thick] (0,1)--(1.5,-.5)--(3,.25)--(4.5,-.125)--(6,.0625)--(7.5,-.03125);
\draw[red!75!black,thick] (0,1)--(2.2,-1.2)--(4.4,1.44)--(6.6,-1.728);
\foreach \t/\y in {0/1,2.2/-1.2,4.4/1.44,6.6/-1.728}
\fill[red!75!black] (\t,\y) circle (1.7pt);
\foreach \t in {2,4,6,8} \node[below] at (\t,0) {\scriptsize $\t$};
\node[left] at (0,1) {$1$}; \node[left] at (0,-1) {$-1$};
\end{tikzpicture}
\end{column}
\begin{column}{.38\textwidth}
Задача: $x'=-x$, $x(0)=1$.
\medskip

{\color{blue!75!black}Синяя кривая}: точное решение.
\medskip

{\color{teal!75!black}Зелёная ломаная}: $h=1.5$, затухает со сменой знака.
\medskip

{\color{red!75!black}Красная ломаная}: $h=2.2$, модуль растёт.
\end{column}
\end{columns}
\end{frame}

% 33
\begin{frame}{Проверьте себя: что именно даёт каждый метод?}
\begin{enumerate}
\item Какое решение теряется при делении уравнения $x'=x^2$ на $x^2$?
\item Может ли непостоянное решение $x'=q(x)$, $q\in C^1$,
быть периодическим?
\item Достигает ли логистическое решение уровня $K$ за конечное время,
если $0<x_0<K$?
\item Почему интегрирующий множитель не теряет решения?
\item Почему $O(h^2)$ за один шаг не означает $O(h^2)$ на всём отрезке?
\item Для $x'=-10x$ шаг $h=0.3$ даёт затухание или рост?
\end{enumerate}
\begin{block}{К следующей лекции}
Мы получили методы и примеры. Теперь нужны общие условия
существования и единственности решения задачи Коши.
\end{block}
\end{frame}

% 34
\begin{frame}{Упражнения A: техника и проверка}
\begin{enumerate}
\item Решите $x'=2t(1+x^2)$, $x(0)=0$.
Укажите максимальный интервал и проверьте решение.
\item Найдите все решения $x'=3t(x-2)^2$.
Не забудьте постоянные решения и интервалы определения.
\item Постройте фазовые прямые для $x'=x^2$,
$x'=x^3$ и $x'=-x^3$. Что показывает сравнение при $q'(0)=0$?
\item Решите $x'+2x=4$, $x(0)=1$.
Найдите предел при $t\to+\infty$ двумя способами.
\item Решите $x'+(2/t)x=t$, $x(1)=0$.
На каком максимальном интервале определено это решение исходного уравнения?
\end{enumerate}
\end{frame}

% 35
\begin{frame}{Упражнения B: модель, геометрия, вычисление}
\small
\begin{enumerate}
\setcounter{enumi}{5}
\item В логистической модели $K=1000$, $x(0)=100$, $x(2)=250$.
Время измеряется в днях. Найдите $r$ и время достижения $x=500$.
\item Для $x'=-2x$, $x(0)=1$ выполните два шага Эйлера при $h=0.25$.
Сравните результат в $t=0.5$ с точным значением.
При каких шагах сеточное решение затухает? При каких остаётся положительным?
\item Для $x'=x(1-x)$, $x(0)=2$, шаг $h=2$ даёт $X_1=-2$.
Объясните, почему это противоречит качественному поведению точного решения.
Какое условие на $h$ обеспечивает хотя бы $X_1\ge0$?
\item Для $x'=x^2$, $x(0)=1$ метод Эйлера с шагом $1/N$
выдаёт конечное число $X_N$ в момент $t=1$.
Почему это не доказывает существования $x(1)$?
\end{enumerate}
\end{frame}

% 36
\begin{frame}{Классические ориентиры и дальнейшее чтение}
Материалы MIT OpenCourseWare, курс 18.03SC:
\begin{itemize}
\item \href{https://ocw.mit.edu/courses/18-03sc-differential-equations-fall-2011/pages/unit-i-first-order-differential-equations/basic-de-and-separable-equations/}{Basic DE's and Separable Equations}.
\item \href{https://ocw.mit.edu/courses/18-03sc-differential-equations-fall-2011/pages/unit-i-first-order-differential-equations/first-order-autonomous-differential-equations/}{First Order Autonomous Differential Equations}.
\item \href{https://ocw.mit.edu/courses/18-03sc-differential-equations-fall-2011/pages/unit-i-first-order-differential-equations/solution-by-integrating-factors/}{Solution by Integrating Factors}.
\item \href{https://ocw.mit.edu/courses/18-03sc-differential-equations-fall-2011/pages/unit-i-first-order-differential-equations/numerical-methods/}{Numerical Methods}.
\end{itemize}
\begin{block}{Как пользоваться}
Сопоставляйте вычисления с фазовой прямой и численным экспериментом.
Подробные комментарии, решения упражнений и ссылки на конспекты:
\texttt{lecture2-script.md}.
\end{block}
Эта лекция --- самостоятельное изложение для нашего курса, не перевод MIT.
\end{frame}
\end{document}
